% Complete replacement draft for q5-c168.tex. Insert both shared files first.
\subsection{The terminal pair (2,3) at a root (1,2,v)}
\label{proof:Q5-C168}

Let \(v\ge3\) and let the original inner and outer connector lengths
\(L,M\ge1\) be independent. Either terminal direction may be fixed:
\begin{equation}\label{eq:odd168:directions}
 F=\bar f_2W_3\quad\hbox{or}\quad F=h_2f_3,
 \qquad
 P_{v,L}=\begin{cases}h_2f_v,&L=2,\\
                       \bar f_2W_v,&L=1\text{ or }L\ge3.
          \end{cases}
\end{equation}
Use the outward unit arm \(p\), direct inner edges from terminal to
root and outer edges from root to receiver, and put
\begin{equation}\label{eq:odd168:statement}
 H=pP_{v,L}T^LF.
 \qquad
 \E H^4\le e^{-25d/4},\qquad
 I_q(T^MH)^3\le e^{-3d/2}.
\end{equation}
The directions depend only on the original lengths. Arbitrary genuine
equal pairs and subdivision vertices are allowed. We use the notation
and path inputs of Section~\ref{sec:odd-common-calculus}.

\begin{proof}
The arm sums for both terminal directions are \(2-d+R\), with
\(R=-B^2u\):
\[
 \bar f_2=1+u+Bu,\quad W_3=1-d-u-Bu-B^2u,
 \quad h_2=1+u-Bu,\quad f_3=1-d-u+Bu-B^2u.
\]
Thus \(F\le e^{R-d}\). The row bound gives
\(|R|\le\sqrt{\delta_2}\le\sqrt{13/96}<3/8\).
On \([0,3/8]\), the Taylor remainder coefficient is at most
\[
 \tfrac12+\tfrac t6+\frac{t^2}{24(1-t/5)}
 \le\tfrac{1347}{2368}<\tfrac47;
\]
on negative arguments it is at most \(1/2\).
Use the positive majorant \eqref{eq:oc:fourth-duhamel}, with the
coefficient \(c=\eta=8/7\) permitted at \(t=3/8\), and data
\begin{equation}\label{eq:odd168:terminal-data}
\begin{gathered}
 \Gamma=d,\quad\eta=8/7,\quad\sigma=\delta_3,
 \quad D_{2k}=0,\quad D_{2k-1}=(-1)^k\delta_{k+1}\ (k\ge1),\\
 \delta=8\delta_3/7\le169d/896<d/5,
 \quad\delta\le13\delta_2/28,
 \quad\epsilon=\delta-\delta_2<0.
\end{gathered}
\end{equation}
Write \(m=v-2\), \(G_m=E+zC\), and \(\beta=\int C\,d\nu\).
The standard and alternative root losses are \(\ell_s,\ell_a\)
from Lemma~\ref{lem:odd-common:root-pairs}.

For the three short lengths, combine their positive majorants with
the indicated root pair and the actual \(p\) curvature. The resulting
parameters \((Q_L,c_L,S_L)\), where \(S_L=a_L+\E BQ_L\), are
\begin{equation}\label{eq:odd168:short-parameters}
\begin{aligned}
 Q_1&=[x(1+C)-zE]u,\qquad c_1=\beta,\\
 S_1&=-d+\delta-\delta_2-\ell_s,\\[2pt]
 Q_2&=[1-xC+z(x-E)]u,\qquad c_2=\delta+\beta,\\
 S_2&=-2d-\delta_2+\delta-\ell_a,\\[2pt]
 Q_3&=[1+xC-x^2+\delta+z(1-E)]u,\qquad c_3=\beta+\epsilon,\\
 S_3&=-3d-\delta_2+\delta+\delta_3-\delta d-\ell_s.
\end{aligned}
\end{equation}
Here \(z=i\sqrt x\); the coordinate notation defines vectors by
functional calculus, not pointwise functions of the original labels.
In every row \eqref{eq:odd-common:curvature} gives
\(\log\E H^4\le4S_L+(17/4)\|Q_L\|^2+(10/3)c_L^2\).

\emph{Length one.}
For \(v=3\), \((E,C,\beta)=(1,0,0)\), giving
\[
 \log\E H^4\le-8d+\tfrac14\delta_2+\tfrac{247}{28}\delta_3
 <-8d+\tfrac{529}{336}d<-25d/4.
\]
For \(v\ge4\), pay \(\beta^2\le d\int C^2\,d\nu\) first.
The complete frequency integrands, for even and odd \(m\), are
\begin{align*}
 F_e={}&-4-4x+\tfrac{32}{7}x^2-(8+4x)C
       +\tfrac{17}{4}[x^2(1+C)^2+xC^2]+\tfrac{10}{9}C^2,\\
 F_o={}&-8+\tfrac x4+\tfrac{247}{28}x^2-4C
       +[\tfrac{17}{4}x^2(1+x)+\tfrac{10}{9}]C^2.
\end{align*}
They are convex in \(C\in[1-x,1]\). All four endpoint polynomials
and their upper bounds on \([0,\theta]\) are
\begin{center}\small
\begin{tabular}{@{}p{.13\linewidth}p{.62\linewidth}p{.19\linewidth}@{}}
\toprule
\((m,C)\)&Endpoint polynomial&Upper bound\\ \midrule
Even, \(1\)&\(-\tfrac{98}{9}-\tfrac{15}{4}x+\tfrac{151}{7}x^2\)
       &\(-\tfrac{919}{126}\)\\[3pt]
Even, \(1-x\)&\(\begin{aligned}[t]
 &-\tfrac{98}{9}+\tfrac{73}{36}x+\tfrac{2291}{126}x^2\\
 &\quad-\tfrac{51}{4}x^3+\tfrac{17}{4}x^4
 \end{aligned}\)&\(-\tfrac{10349}{1512}\)\\[3pt]
Odd, \(1\)&\(-\tfrac{98}{9}+\tfrac x4+\tfrac{183}{14}x^2+\tfrac{17}{4}x^3\)
       &\(-\tfrac{8297}{1008}\)\\[3pt]
Odd, \(1-x\)&\(\begin{aligned}[t]
 &-\tfrac{98}{9}+\tfrac{73}{36}x+\tfrac{1787}{126}x^2\\
 &\quad-\tfrac{17}{4}x^3-\tfrac{17}{4}x^4+\tfrac{17}{4}x^5
 \end{aligned}\)&\(-\tfrac{11357}{1512}\)\\
\bottomrule
\end{tabular}
\end{center}
The bounds use \(x\le1/2\), \(x^2<1/6\), \(x^3<1/12\).
In the second and fourth rows the final displayed tails are
nonpositive. Every bound is below \(-25/4\).

\emph{Length two.}
For \(v\ge4\), the actual alternative root gives
\[
 \|Q_2\|^2\le\int[(1-x+x^2)^2+x(1-x)^2]\,d\nu
 =\int[1-x+x^2-x^3+x^4]\,d\nu\le d.
\]
Indeed the last polynomial minus one is
\(-x(1-x)(1+x^2)\le0\).
Moreover \(0\le c_2\le6d/5\) and
\(4\delta-4\delta_2\le-15\delta_2/7\). Therefore
\[
 \log\E H^4\le-\tfrac{1929}{256}d-\tfrac{15}{7}\delta_2
                    +\tfrac{24}{5}d^2
 \le-(\tfrac{1929}{256}-\tfrac{31}{35})d<-13d/2.
\]
For \(v=3\), keep the exact norm
\(\|Q_2\|^2=d+\delta_2-2\delta_3+\delta_4\) and \(c_2=\delta\):
\[
 \log\E H^4\le-\tfrac{31}{4}d+\tfrac14\delta_2
        -\tfrac{55}{14}\delta_3+\tfrac{17}{4}\delta_4
        +\tfrac{10}{3}\delta^2<-\tfrac{303}{40}d.
\]
Here \(\delta_4\le\theta\delta_3\) makes the middle pair nonpositive,
and the remaining costs are less than \(d/8+d/20\).

\emph{Length three.}
For \(v\ge4\), \(0\le c_3\le d\); for \(v=3\),
\(c_3=\epsilon<0\) and \(|c_3|\le\delta_2\le d\).
Removing \(\delta u\) only for the norm comparison gives
\[
 \|Q_3-\delta u\|^2\le
       \int[(1+x-x^2)^2+x^3]\,d\nu
       \le d+2\delta_2-\delta_3.
\]
Restoring its square and cross term adds at most
\((5/2)\delta d+\delta^2d\le77d^2/150\).
Drop the favorable \(-4\delta d\) from the scalar term. The full bound is
\[
 \log\E H^4\le-\tfrac{47}{4}d+\tfrac92\delta_2
                 +\tfrac{121}{28}\delta_3+\tfrac{1103}{200}d^2
 <-\tfrac{27}{4}d.
\]
The three positive terms are less than \(9d/4,3d/4,2d\),
respectively. This includes the negative constant at \(v=3\).

\phantomsection\label{part:C168.8}\emph{Every longer connector.}
For \(\gamma_j=\beta_j-D_{j-1}\), the finite moments in
\eqref{eq:odd168:terminal-data} give
\begin{equation}\label{eq:odd168:gamma}
 \gamma_3=d,\quad\gamma_4=d-\delta_3\ge d/2,
 \qquad\gamma_j\ge\kappa d-\delta_4
       >\tfrac{117}{224}d>d/2\quad(j\ge5).
\end{equation}
In the last inequality \(\delta_4\le\theta^3d<d/14\).
Apply \eqref{eq:odd-common:long-parameters} with \(h=1/2\),
\(D_1<0\), \(n=L-1\ge3\), and \(R=-B^2u\).
The retained \(\gamma_L\bar f_0\) and
\(\gamma_{L-1}\bar f_1\) give a terminal scalar-plus-mean bound
\begin{equation}\label{eq:odd168:long-mean}
 S_t\le-\tfrac52d+\tfrac{16}{7}\delta_3
       \le-\tfrac{951}{448}d<-\tfrac{21}{10}d.
\end{equation}
The next mean uses \(\gamma_{L+1}\), covered by the same finite bounds.
Adding the standard root leaves the curvature parameters
\begin{equation}\label{eq:odd168:long-parameters}
 Q=(G_n(B)-BG_m(B)-B^{n+2})u,
 \qquad c=\beta-d/2,\qquad S\le S_t-\ell_s.
\end{equation}
Thus \(|c|\le d/2\), whose curvature cost is at most \(5d/18\).
Also \(\|B^{n+2}u\|<11\sqrt d/100\), since
\(\theta^5<121/10000\).
For \(v\ge4\), Lemma~\ref{lem:odd-common:joint-paths} gives
\(\|Q\|<131\sqrt d/100\), hence \(\|Q\|^2<7d/4\), while
\(S<-7d/2\). Consequently
\[
 \log\E H^4<-14d+\tfrac{119}{16}d+\tfrac5{18}d
        =-\tfrac{905}{144}d<-25d/4.
\]
For \(v=3\), the stronger \(m=1\) bound gives
\(\|Q\|<111\sqrt d/100\), \(\|Q\|^2<5d/4\),
\(S<-31d/10\), and hence
\(\log\E H^4<-4903d/720<-25d/4\).

All \(L\ge1\) are now covered. The uniform cost \(11\) in
\eqref{eq:odd-common:outlet} gives
\(I_q(T^MH)^3\le e^{[11-2(25/4)]d}=e^{-3d/2}\)
for every independent \(M\ge1\).
Genuine equal pairs are restored by positivity and their pointwise
bound in \eqref{eq:odd-common:inputs}. Complete reversal applies the
same proof to \(-K\), exchanging the corresponding directions and
weights. The \(d=0\) boundary follows by leaf integration.
\end{proof}
